% ============================================================
% 分片：s05a-wps-toric  对应原文 PDF 第 49--57 页（印刷页码 49--57）
% 负责章节：§5 GENERALIZATIONS 引言 + 5.1 Hypersurfaces In Weighted Projective Space
%           + 5.2 Hypersurfaces In Toric Varieties
% 译者：t04-s3a   审校状态：初稿
% 审校：t07-s34 已完成逐式核对（公式修正 0 处；批注表述修正 2 处）
% ============================================================
% 边界说明（依 Lead 指示）：第 49 页首行 "…holomorphic function on different complex
% manifolds." 是 §4 末句的跨页续文，已由 s04-geometry 分片按 SPEC §2 补全；本分片
% 自第 49 页 "For instance, in the case treated in §3…" 起笔。因此第 49 页 §5 标题
% 之前的两段（§4 的收尾：Z 与 Z' 上的共同超势）仍由本分片翻译。第 57 页末句
% "Batyrev [40] gave a kind…" 跨至第 58 页，其余部分归 s05b 分片。
% 原文用直立体 d 表示微分，故微分记号统一写作 \mathrm{d}。
% 原文 projective space 排为粗体 CP/WCP，本分片保留为 \mathbf{CP}、\mathbf{WCP}。
% 原文第 56 页脚注（⋆）以 \footnote{} 保留，编号由 LaTeX 重排。

% ==== 原文 PDF 第 49 页（印刷页 49）====
\label{op:49}
% [边界: 本页首行属 §4 末句，已由 s04-geometry 收走；本分片自此段的 "For instance…" 起笔。]

例如，在 §3 所处理的情形中，底层的超势为 $W=pG(s_i)$，其中 $G$ 是 $n$ 次齐次多项式。
对 $r>0$，$W$ 限制并下降到辛商 $Z$ 上，成为全纯函数 $\widehat{W}$。积掉有质量规范多重态
之后的有效理论，是以 $Z$ 为靶空间、以 $\widehat{W}$ 为超势的理论；利用 §3 的（本质上相同的）
计算，它在低能下约化为 $\mathbf{CP}^{n-1}\subset Z$ 中 Calabi--Yau 超曲面 $G=0$ 的 $\sigma$ 模型。

另一方面，对 $r<0$，$W$ 限制并下降到 $Z'=\mathbb{C}^n/\mathbb{Z}_n$ 上的全纯函数 $G(s_i)$。
这就是熟悉的 Landau--Ginzburg 轨形的超势。

这样一来，$Z$ 与 $Z'$ 之间的双有理等价就隐藏在 Calabi--Yau/Landau--Ginzburg 对应的背后。
基本的关系是 $Z$ 上的量子场论与 $Z'$ 上的量子场论之间的关系；而 Calabi--Yau/Landau--Ginzburg
对应，则是在存在某个共同超势的情况下考察这一关系时出现的。

\section{推广}

\begin{keytakeaway}
\textbf{核心物理图像：}§3 的规范线性 $\sigma$ 模型构造可以系统推广：只要把环境空间从
$\mathbf{CP}^{n-1}$ 换成加权射影空间 $\mathbf{WCP}^{n-1}_{q_1,\dots,q_n}$、射影空间之积或更一般的
环面簇，并相应调整 $U(1)^d$ 的规范荷与超势的齐次次数，Calabi--Yau/Landau--Ginzburg
对应便逐字延续。改变 Fayet--Iliopoulos 参数（$r$ 或 $r_j$）使同一模型在不同相之间移动；
新出现的相可以是 $\sigma$ 模型与 Landau--Ginzburg 模型的“混合”，相变还可以伴随靶空间
（时空）拓扑的改变。

\par\textbf{数学工具：}加权射影空间是 $\mathbb{C}^{n+1}$ 对按权 $q_i$ 作用的 $U(1)$ 的辛商；
$U(1)^d$ 的 $D$ 平坦商给出环面簇 $Z$（复维 $N-d$，带 $(\mathbb{C}^*)^{N-d}$ 的环面作用），
超曲面 $X=\{G=0\}\subset Z$ 即环面超曲面，横截性保证 $X$ 在 $Z$ 的奇点之外光滑；
双次数 $(n,m)$ 超曲面的伴随公式给出 Calabi--Yau 条件。

\par\textbf{主要结论：}(1) 当 $q_i$ 互素、$G$ 的荷为 $q=\sum_i q_i$ 时，$r\gg0$ 给出
$\mathbf{WCP}^{n-1}_{q_1,\dots,q_n}$ 中的 Calabi--Yau 超曲面，$r\ll0$ 给出 $\mathbb{Z}_q$ 轨形的
Landau--Ginzburg 模型；(2) $\mathbf{CP}^{n-1}\times\mathbf{CP}^{m-1}$ 模型有三个相，
相 II 与相 III 是 $\sigma$ 模型与 Landau--Ginzburg 轨形的混合，相边界上得到规范
Landau--Ginzburg 模型；(3) 一般环面簇的情形囊括上述一切，并预言相变引起的拓扑改变
（§5.5）以及 Batyrev 构造的场论背景。
\end{keytakeaway}

本节的任务是描述 Calabi--Yau/Landau--Ginzburg 对应的一些推广。我们首先考虑若干较为
直接的推广，涉及：(i) 加权射影空间中的超曲面；(ii) 射影空间之积以及更一般的环面簇中的
超曲面；(iii) 格拉斯曼流形中的超曲面；(iv) 上述任一空间中超曲面的交。每一种情形下的
讨论都将相当简短，因为 §3 中的大部分讨论只需作显然的修改即可照搬。最后，在 §5.5 中
我们将讨论不同拓扑的流形之间的跃迁如何在这一框架中出现。§4 的框架是有用的背景，
尤其对于 §5.5。

% ==== 原文 PDF 第 50 页（印刷页 50）====
\label{op:50}

\subsection{加权射影空间中的超曲面}

对 §3 中已有结果所作的最明显的轻微推广如下。考虑一个 $U(1)$ 规范理论，它含 $n$ 个
手征超场 $S_i$（$i=1,\dots,n$），荷为 $q_i$；再含一个手征超场 $P$，荷为 $-\sum_i q_i$。
通过对 $U(1)$ 生成元重新标度，可以假定 $q_i$ 互素。
% [Tutor Note: “对 $U(1)$ 生成元重新标度”即把同一 $U(1)$ 的生成元乘以非零常数，等价于把全部荷同乘一个常数；这不改变规范等价类，只重新定义规范耦合 $e$ 与 FI 参数 $r$。故可把荷的公因子归一化，取 $\gcd(q_1,\dots,q_n)=1$ 不失一般性。]

上述荷的选取保证了 $\sum_i Q_i=0$；正如我们在 §3 中所见，这一条件导致无反常的
左行与右行 $R$ 对称性，从而在低能下得到 Calabi--Yau 流形。在随后的例子中，我们将
始终以类似方式保证对每一个 $U(1)$ 荷都有 $\sum_i Q_i=0$，而不在每种情形下逐一指出这一点。

我们把超势取为
\begin{equation}
W=PG(S_i),
\tag{5.1}
\end{equation}
其中 $G$ 是荷为 $q=\sum_i q_i$ 的多项式，并且是横截的，意即方程
\begin{equation}
\frac{\partial G}{\partial S_1}=\dots=\frac{\partial G}{\partial S_n}=0
\tag{5.2}
\end{equation}
仅在 $S_1=\dots=S_n=0$ 时成立。（这样的 $G$ 的存在性是对 $q_i$ 的严厉限制，[37] 中对此作了
分析；但只要存在任何一个荷为 $q$ 的 $G$ 满足该条件，一般的 $G$ 就满足。）于是，拉格朗日量中
含手征超场的部分为
\begin{equation}
\tag{5.3}
\begin{aligned}
&\int \mathrm{d}^2y\,\mathrm{d}^4\theta\Bigl(\sum_i\bar S_ie^{2V}S_i+\bar Pe^{-2qV}P\Bigr)
-\int \mathrm{d}^2y\,\mathrm{d}^2\theta\,PG(S_i)\\
&\qquad-\int \mathrm{d}^2y\,\mathrm{d}^2\bar\theta\,\bar P\,\overline{G}(\bar S_i).
\end{aligned}
\end{equation}
% [Tutor Note: (5.3) 只写出含手征超场的部分，规范部分（规范动能加 Fayet--Iliopoulos 项）未写出。三项依次为：手征超场的动能（$e^{2V}$ 因子保证规范不变性，$P$ 的指数 $-2qV$ 反映其荷为 $-q$）、超势的 $F$ 项及其厄米共轭。注意 $\sum_iQ_i=0$ 之所以能同时保证左右行 $R$ 对称性无反常，是因为它使 $U(1)$ 的规范反常与 $R$ 反常同时抵消；这正是低能极限成为 Calabi--Yau $\sigma$ 模型的原因。]

势能与以前大致相同
\begin{equation}
\tag{5.4}
U(A_i,\sigma)=\frac{1}{2e^2}D^2+\sum_\alpha|F_\alpha|^2
+2\bar\sigma\sigma\Bigl(\sum_iq_i^2|s_i|^2+q^2|p|^2\Bigr),
\end{equation}
% [存疑: (5.4) 左端印作 $U(A_i,\sigma)$（图像确认），但右端及后文（如 (5.11)）均以玻色分量 $s_i$ 记标量场；$A_i$ 疑为早期记号之遗留。照录原文，此处存疑待考。]

% ==== 原文 PDF 第 51 页（印刷页 51）====
\label{op:51}

其中
\begin{equation}
D=-e^2\Bigl(\sum_iq_i|s_i|^2-q|p|^2-r\Bigr)
\tag{5.5}
\end{equation}
并且
\begin{equation}
\sum_\alpha|F_\alpha|^2=|G|^2+|p|^2\sum_i\left|\frac{\partial G}{\partial s_i}\right|^2.
\tag{5.6}
\end{equation}

由于我们将不得不反复写出类似的公式，先作几点约定。符号 $\sum_\alpha$ 始终表示对全部
手征超场 $\Phi_\alpha$ 求和，其辅助场为 $F_\alpha$。势中涉及 $\sigma$ 的项总是具有给 $\sigma$
以质量、迫使其真空期望值消失的效果（远离相边界处；相边界附近的行为总如 §3.2 所述）。
知道了这一点，就可以在讨论中丢弃 $\sigma$，只在 $\sigma=0$ 处写出势。
% [Tutor Note: $\sigma$ 的质量平方直接来自 (5.4) 的最后一项：$2\bar\sigma\sigma(\sum_iq_i^2|s_i|^2+q^2|p|^2)$ 中 $|\sigma|^2$ 的系数即 $\sigma$ 的有效质量平方（与 §3.2 的讨论一致，因子 2 取决于规范多重态的归一化约定）。在 $\langle s_i\rangle$ 或 $\langle p\rangle$ 非零处 $\sigma$ 有质量，可被积掉；只有在相边界上该系数可能为零，$\sigma$ 才重新变成无质量场，此即 §3.2 所说的“奇点”附近的物理。]

分析完全照 §3.1 进行。确定基态结构的一步是令 $D=0$ 并对规范群 $U(1)$ 取商（用 §4 的语言，
即作 $\mathbb{C}^{n+1}$ 对 $U(1)$ 的辛商）。若 $r\gg0$，$D=0$ 的轨迹对 $U(1)$ 取商后就是
加权射影空间 $\mathbf{WCP}^{n-1}_{q_1,\dots,q_n}$ 的一份拷贝，其凯勒类正比于 $r$。低能理论是以
该加权射影空间中的超曲面 $G=0$ 为靶空间的 $\sigma$ 模型。由于底层的 $R$ 不变性，该超曲面
是 Calabi--Yau 流形。$q_i$ 互素的假定意味着（在加权射影空间的奇点之外）规范对称性被
完全破缺。
% [Tutor Note: 加权射影空间的标准定义是 $(\mathbb{C}^n\setminus\{0\})/\mathbb{C}^*$，其中 $\lambda\in\mathbb{C}^*$ 按 $s_i\to\lambda^{q_i}s_i$ 作用；它与这里的辛商一致（矩映射水平为 $r$）。其奇点位于 $U(1)$ 稳定化子非平凡的层上：若只剩 $i\in J$ 的坐标非零，稳定化子为 $\mu_d$（$d=\gcd_{i\in J}q_i$；单坐标时 $d=q_i$），商空间局部是 $\mathbb{C}^{n-|J|}/\mu_d$（$|J|$ 为该层中非零坐标的个数）型商奇点。$q_i$ 互素只保证“一般点”（$s_i$ 全不为零、$p=0$）处稳定化子平凡，即规范对称性完全破缺，并不消除低维层上的奇点；弦论上这些奇点由世界面轨形（及其离散挠率自由度）处理。]

另一方面，对 $r\ll0$，场 $p$ 获得真空期望值，把规范群破缺为 $\mathbb{Z}_q$。在规范变换意义下
存在唯一真空：$s_i=0$，$p=\sqrt{-r/q}$。低能理论是超势为 $G(s_1,\dots,s_n)$ 的
$\mathbb{Z}_q$ 轨形。我们就这样重新得到了加权射影空间中 Calabi--Yau 超曲面与
Landau--Ginzburg 轨形之间的标准对应。
% [Tutor Note: 残余群阶数的来源：规范变换为 $s_i\to e^{iq_i\lambda}s_i$、$p\to e^{-iq\lambda}p$。要保持 $\langle p\rangle\neq0$ 不变须 $e^{-iq\lambda}=1$，即 $e^{i\lambda}$ 为 $q$ 次单位根，故 $E\cong\mathbb{Z}_q$；作用在 $s_i$ 上的权为 $q_i$。它与 §3.1 中 $\mathbb{Z}_n$ 的情形一致（那里 $q=n$）。]

% ==== 原文 PDF 第 52 页（印刷页 52）====
\label{op:52}

\subsection{环面簇中的超曲面}

现在要考察一种将使我们发现新东西、而不只是重新得到已知结果的情形。为避免记号变得
累赘，先考虑射影空间之积中超曲面的一个非常特殊的情形，比如说
$\mathbf{CP}^{n-1}\times\mathbf{CP}^{m-1}$。

首先，为了描述 $\mathbf{CP}^{n-1}\times\mathbf{CP}^{m-1}$，考虑一个带两个矢量超场（比如说
$V_1$ 与 $V_2$）的 $U(1)\times U(1)$ 规范理论，规范耦合为 $e_1$ 与 $e_2$。因此有两个独立的
Fayet--Iliopoulos 项，其分量形式为
\begin{equation}
L_D=-\int\mathrm{d}^2y\,(r_1D_1+r_2D_2).
\tag{5.7}
\end{equation}

我们引入荷为 $(1,0)$ 的超场 $S_i$（$i=1,\dots,n$）、荷为 $(0,1)$ 的超场 $T_j$
（$j=1,\dots,m$），并再引入一个荷为 $(-n,-m)$ 的超场 $P$。于是这些场的动能为
\begin{equation}
\tag{5.8}
\begin{aligned}
&\int\mathrm{d}^2y\,\mathrm{d}^4\theta\Bigl(\sum_i\bar S_ie^{2V_1}S_i
+\sum_j\bar T_je^{2V_2}T_j\\
&\qquad+\bar Pe^{-2nV_1-2mV_2}P\Bigr).
\end{aligned}
\end{equation}

超势取为
\begin{equation}
W=PG(S_i,T_j),
\tag{5.9}
\end{equation}
其中 $G$ 是一个在 $S_i$ 中 $n$ 次齐次、在 $T_j$ 中 $m$ 次齐次的多项式，并且是横截的，
意即方程
\begin{equation}
0=\frac{\partial G}{\partial S_i}=\frac{\partial G}{\partial T_j}
\tag{5.10}
\end{equation}
除“$S_i$ 全为零”或“$T_j$ 全为零”的情形外无解。该理论的势能

% ==== 原文 PDF 第 53 页（印刷页 53）====
\label{op:53}
% [跨页接续: 第 52/53 页；"The potential energy / of this theory is (at σ = 0)" 跨页，按语义合译。]
（在 $\sigma=0$ 处）为
\begin{equation}
\tag{5.11}
\begin{aligned}
U(s_i,t_j)&=\frac{e_1^2}{2}\Bigl(\sum_i|s_i|^2-n|p|^2-r_1\Bigr)^2
+\frac{e_2^2}{2}\Bigl(\sum_j|t_j|^2-m|p|^2-r_2\Bigr)^2\\
&\qquad+\sum_\alpha|F_\alpha|^2,
\end{aligned}
\end{equation}
其中
\begin{equation}
\tag{5.12}
\sum_\alpha|F_\alpha|^2=|G|^2+|p|^2\Bigl(\sum_i\left|\frac{\partial G}{\partial s_i}\right|^2
+\sum_j\left|\frac{\partial G}{\partial t_j}\right|^2\Bigr).
\end{equation}

现在分析真空结构。$G$ 的横截性意味着 (5.12) 消失当且仅当以下三者之一成立：(i) $p=0$ 且
$G=0$；(ii) $s_i=0$（$i=1,\dots,n$），其余场取适当值；(iii) $t_j=0$（$j=1,\dots,m$），
其余场取适当值。这三种选择对应理论的三个相，我们分别称为相 I、相 II 与相 III。

实现哪一个相取决于 $r_i$ 的取值。令 (5.11) 为零即可发现：相 I 在 $r_1,r_2\ge0$ 时实现；
相 II 在 $r_1\le0$、$nr_2-mr_1\ge0$ 时实现；相 III 在 $r_2\le0$、$nr_2-mr_1\le0$ 时实现。
因此相边界是 $r_1$-$r_2$ 平面中的三条半直线：$r_1=0$、$r_2\ge0$；$r_2=0$、$r_1\ge0$；
以及 $r_1,r_2\le0$、$nr_2-mr_1=0$。这些相边界把平面分成相 I、相 II 或相 III 的区域。
% [Tutor Note: 相图可逐条从 $U=0$ 解出：两个 $D$ 条件给出 $\sum_i|s_i|^2-n|p|^2=r_1$、$\sum_j|t_j|^2-m|p|^2=r_2$，$F$ 项给出 $G=0$ 与 $p\,\partial G/\partial s_i=p\,\partial G/\partial t_j=0$。分支 (i)（$p=0$）要求 $r_1,r_2\ge0$；分支 (ii)（$s_i=0$，$p\neq0$）由第一式得 $|p|^2=-r_1/n\ge0$，代入第二式得 $\sum_j|t_j|^2=r_2-mr_1/n\ge0$，即 $r_1\le0$、$nr_2-mr_1\ge0$；分支 (iii) 对称地给出 $r_2\le0$、$nr_2-mr_1\le0$。边界线正是相区交界处：在 $r_1=0$ 或 $r_2=0$ 上，相 I 与混合相在 $s=0$ 或 $t=0$ 处相接。]

相 I 是熟悉的 Calabi--Yau 相。在这个相中，$s_i$ 与 $t_j$ 的期望值把规范群完全破缺。
真空中 $p=0$，$\sum_i|s_i|^2=r_1$，$\sum_j|t_j|^2=r_2$。对规范群 $U(1)\times U(1)$ 取商后，
$s_i$ 与 $t_j$ 确定 $\mathbf{CP}^{n-1}\times\mathbf{CP}^{m-1}$ 中的一个点，其凯勒类正比于
$r_1,r_2$。由于 (5.12) 的消失还要求 $G=0$（且所有不沿 $G=0$ 的解空间切向的模式都有质量），
低能理论是 $\mathbf{CP}^{n-1}\times\mathbf{CP}^{m-1}$ 中超曲面 $G=0$ 的 $\sigma$ 模型。由于 $G$
的双次数为 $(n,m)$，这实际上是一个 Calabi--Yau 超曲面。
% [Tutor Note: 双次数 $(n,m)$ 给出 Calabi--Yau 条件可由伴随公式核对：$K_{\mathbf{CP}^{n-1}\times\mathbf{CP}^{m-1}}=\mathcal{O}(-n,-m)$，而 $X$ 是 $\mathcal{O}(n,m)$ 的零点集，故 $K_X=K_Z\otimes\mathcal{O}(n,m)\big|_X=\mathcal{O}_X$，即 $c_1(X)=0$。这与 §3 中“$P$ 的荷决定 $G$ 的次数、无反常条件迫使次数等于环境空间维数”的机制完全平行：此处 $P$ 的荷 $(-n,-m)$ 同时锁定了两个次数。]

现在考察相 II。在这个区域，$p$ 与 $t_j$ 有真空期望值，把规范群自发破缺为 $\mathbb{Z}_n$。
这些期望值

% ==== 原文 PDF 第 54 页（印刷页 54）====
\label{op:54}
% [跨页接续: 第 53/54 页；"The expectation values / obey (5.13)" 跨页，按语义合译。]
满足
\begin{equation}
|p|=\sqrt{-r_1/n},\qquad \sum_j|t_j|^2=r_2-\frac{mr_1}{n}.
\tag{5.13}
\end{equation}
通过一个规范变换，可以假定 $p>0$。这留下残余的规范不变性
\begin{equation}
t_j\to\exp(i\theta)\cdot t_j,\qquad s_i\to\exp(-i(m/n)\theta)\cdot s_i.
\tag{5.14}
\end{equation}
在施加 (5.13) 中第二个方程与规范不变性 (5.14) 之后，$t_j$ 的取值确定
$\mathbf{CP}^{m-1}$ 中的一个点。由于 $s_i$ 在该规范变换下按 (5.14) 变换，应当认为 $s_i$
并不定义 $\mathbb{C}^n$ 中的一个点，而是定义 $\mathbf{CP}^{m-1}$ 上一个 $n$ 维复矢量丛 $Y$
的纤维中的一个点。（若 $m/n$ 不是整数，这就不是普通的矢量丛，而是 V-丛（V-bundle），
即轨形的矢量丛类比。）把 $p$ 取为其真空期望值后，存在一个由 $Y$ 上的全纯函数
$W_{eff}=\sqrt{-r_1/n}\cdot G(s_i,t_j)$ 给出的有效超势。因此相 II 是 $t$ 方向上的
$\mathbf{CP}^{m-1}$ $\sigma$ 模型与 $s$ 方向上的 Landau--Ginzburg 轨形的一种奇特混合。
这是我们第一次遇到这样的混合，但我们将看到，各种类型的混合比比皆是。
% [Tutor Note: 为何 $s_i$ 只定义丛 $Y$ 的纤维坐标：残余规范变换 (5.14) 只把 $s_i$ 整体乘以 $e^{-i(m/n)\theta}$，故在底空间 $\mathbf{CP}^{m-1}$ 上 $s$ 不是函数而是截面（纤维坐标）——每个 $t$ 给出底空间的一个点，$s$ 则生活在附加的 $\mathbb{C}^n$ 纤维中。当 $m/n\in\mathbb{Z}$ 时过渡函数良定义，$Y$ 为普通全纯矢量丛；否则绕闭路一周只能回到相差一个相位 $e^{-2\pi i m/n}$ 的纤维，这正是 V-丛（相当于把世界面轨形的离散挠率/多值性编码进丛的过渡函数）的定义。$W_{eff}$ 作为 $Y$ 上的函数在纤维方向是全纯的（$G$ 齐次），故定义无歧义。]

相 III 当然可用同样方式处理，它是 $\mathbf{CP}^{n-1}$ $\sigma$ 模型与 Landau--Ginzburg 轨形的
类似混合。

人们或许会问：能否用一个更普通的 Landau--Ginzburg 模型来描述射影空间之积中的
Calabi--Yau 超曲面？我们能接近这一点的最佳做法是走到相 II 与相 III 之间的相边界，
即取 $r_1,r_2<0$ 且 $mr_1-nr_2=0$。此时在规范变换意义下有唯一的经典真空：
$s_i=t_j=0$，$p=\sqrt{-r_1/n}=\sqrt{-r_2/m}$。$p$ 的真空期望值把规范群破缺为
$H=U(1)\times\mathbb{Z}_d$，其中 $d$ 是 $n$ 与 $m$ 的最大公因数。低能理论是规范群为 $H$
的规范 Landau--Ginzburg 模型（gauged Landau--Ginzburg model）。

% ==== 原文 PDF 第 55 页（印刷页 55）====
\label{op:55}

作为一个例子，取 $n=m=3$，$G(S_i,T_j)$ 为双次数 $(3,3)$ 的齐次多项式。在相 II 与相 III
的边界上，未破缺的规范群为 $H=U(1)\times\mathbb{Z}_3$；$U(1)$ 按 $S_i\to\exp(i\theta)S_i$、
$T_j\to\exp(-i\theta)T_j$ 作用，而 $\mathbb{Z}_3$ 用 $1$ 的一个三次方根乘 $S_i$ 而保持 $T_j$
不变。低能理论由有效作用量
\begin{equation}
\tag{5.15}
\begin{aligned}
L_{eff}={}&L_{gauge}+\int\mathrm{d}^2y\,\mathrm{d}^4\theta
\Bigl(\bar S_ie^{2V}S_i+\bar T_je^{-2V}T_j-rV\Bigr)\\
&-\int\mathrm{d}^2y\Bigl(\int\mathrm{d}^2\theta\,\langle p\rangle G+\mathrm{h.c.}\Bigr),
\end{aligned}
\end{equation}
描述，其中参数 $r$ 度量到 II--III 相边界的距离。恰在 $r=0$ 时，模型在原点有孤立的经典真空
且规范不变性未破缺。此时模型在低能下就是一个规范 Landau--Ginzburg 模型。
% [Tutor Note: (5.15) 中已只剩一个矢量超场 $V$：$H=U(1)\times\mathbb{Z}_3$ 的 $U(1)$ 部分由原来两个 $U(1)$ 的对角组合生成，它按 $S_i\to e^{i\theta}S_i$、$T_j\to e^{-i\theta}T_j$ 作用，故 FI 项写成 $\int\mathrm{d}^4\theta(-rV)$（其分量形式为 $-rD$）。$r$ 是 $r_1,r_2$ 在垂直于相边界方向上的线性组合（并已按上式归一化），故 $r=0$ 恰对应 II--III 相边界；此时 $\langle s\rangle=\langle t\rangle=0$、$\langle p\rangle\neq0$，全部规范群未破缺，这个孤立真空由规范 Landau--Ginzburg 模型描述。]

规范 Landau--Ginzburg 模型可以像普通模型那样分析。对普通 Landau--Ginzburg 模型，
未扭扇区对 $(c,c)$ 手征环（即 $B$ 模型的手征环）的贡献，是手征场中多项式 $A(S_i,T_j)$
模去由 $\partial G/\partial S_i$、$\partial G/\partial T_j$ 生成的通常关系所得的环。对
Landau--Ginzburg 轨形，要求 $A$ 在相应的离散对称性下不变。对规范 Landau--Ginzburg 模型，
$A$ 必须是规范不变的。（正如我们在 §3.1 中学到的，轨形不过是具有有限规范群的规范理论的
特例。）在目前的情形下，规范不变性可以这样施加：$S_i$ 与 $T_j$ 中的 $U(1)$ 不变齐次多项式
必须在 $S_i$ 与 $T_j$ 中次数相等，而 $\mathbb{Z}_3$ 不变性意味着该次数必须被 3 整除。
于是手征环由双次数为 $(0,0)$、$(3,3)$、$(6,6)$ 与 $(9,9)$ 的多项式构成。这个模型似乎
不太可能用普通的（不带规范群的）Landau--Ginzburg 模型来描述。若不作连续规范不变性的要求，
$(c,c)$ 手征环将是无限维的。
% [Tutor Note: 手征环的筛法：$H$ 中的 $U(1)$ 按荷 $(+1,-1)$ 作用，故不变单项式 $S^{a}T^{b}$ 须满足 $\sum a_i=\sum b_j$（即双次数 $(k,k)$）；$\mathbb{Z}_3$ 再把 $S_i$ 乘 $\omega$，要求 $k\equiv0\pmod3$。在此约束下原文指出只有 $(0,0),(3,3),(6,6),(9,9)$ 四个双次数留存；更高双次数由 $\partial G/\partial S_i$、$\partial G/\partial T_j$ 生成的关系消去（具体截断依 $G$ 的选取而定，原文未展开）。若完全去掉连续规范不变性的要求，就没有“等次”这一约束，$(c,c)$ 手征环将是无限维的——这正是普通不带规范群的 LG 模型无法刻画该模型的原因。]

\subsubsection*{更一般的环面簇}

一个推广是考虑含 $N$ 个手征超场 $S_1,\dots,S_N$、再含一个手征超场 $P$、规范群为
$H=U(1)^d$ 的模型。在前一

% ==== 原文 PDF 第 56 页（印刷页 56）====
\label{op:56}
% [跨页接续: 第 55/56 页；"In the previous / example" 跨页，按语义合译。]
例中，$N=n+m$ 且 $d=2$。可以把 $S_i$ 看作 $V=\mathbb{C}^N$ 上的线性函数。设第 $j$ 个
$U(1)$ 按 $S_i\to\exp(i\theta q_{ij})S_i$ 作用（$q_{ij}$ 为整数），并按
$P\to\exp(-i\theta q_j)P$ 作用，其中 $q_j=\sum_i q_{ij}$。我们把超势取为
$W=PG(S_1,\dots,S_N)$，其中 $G$ 在第 $j$ 个 $U(1)$ 下的荷为 $q_j$。势能为
\begin{equation}
U(s_1,\dots,s_n,p)=\sum_j\frac{1}{2e_j^2}D_j^2+\sum_\alpha|F_\alpha|^2
\tag{5.16}
\end{equation}
% [存疑: (5.16) 左端写 $U(s_1,\dots,s_n,p)$，但本段引入的是 $N$ 个手征超场 $S_1,\dots,S_N$，且 (5.17) 的求和指标 $i$ 遍及全部 $N$ 个场；此处 $s_n$ 疑应为 $s_N$。照抄原样，待核。]
其中
\begin{equation}
D_j=-e_j^2\Bigl(\sum_iq_{ij}|s_i|^2-q_j|p|^2-r_j\Bigr)
\tag{5.17}
\end{equation}
并且
\begin{equation}
\sum_\alpha|F_\alpha|^2=|G|^2+|p|^2\sum_i\left|\frac{\partial G}{\partial s_i}\right|^2.
\tag{5.18}
\end{equation}

若暂时限制到 $p=0$，则令 $D_j=0$ 并对 $H$ 取商，就得到 $V=\mathbb{C}^N$ 对 $H$ 的所谓
辛商，如 §4 所讨论。这个商记作 $Z$，实维数为 $2(N-d)$，并且显然带有群 $F=U(1)^{N-d}$ 的
作用（即 $s_i$ 的、未被包含在 $H$ 中的那些相位旋转）。事实上（在选定 $r_j$ 之后），$Z$
自然地是凯勒流形——这在 $N=2$ 超对称的形式体系中是显然的；在数学上，$Z$ 上的复结构
可利用辛商与全纯商之间的关系来构造，如 §4 所勾勒。利用 $Z$ 上的复结构，$F$ 在 $Z$ 上的
作用可以解析延拓为 $F_{\mathbb{C}}=(\mathbb{C}^*)^{N-d}$ 的作用。\footnote{作用在
$Z=V//H$ 上的 $F_{\mathbb{C}}$ 作用可以具体构造如下。$V$ 上的 $F$ 作用由变换
$s_i\to\lambda^{e_i}s_i$（$\lambda\in U(1)$，$e_i$ 为某些指数）生成。允许
$\lambda\in\mathbb{C}^*$，就把 $V$ 上的 $F$ 作用延拓为 $F_{\mathbb{C}}$ 作用；由于它与 $H$
交换，它下降到 $V/H_{\mathbb{C}}$ 上的 $F_{\mathbb{C}}$ 作用，而按 §4 所知，后者在一个
稠密开集上与 $V//H$ 相同。}$Z$ 与 $F_{\mathbb{C}}$ 具有相同的复维数 $N-d$，且
$F_{\mathbb{C}}$ 自由作用于 $Z$ 中的一个稠密开集；因此 $Z$ 是所谓的环面簇（toric variety）。
% [Tutor Note: 实维数 $2(N-d)$ 的来历：$V=\mathbb{C}^N$ 实维 $2N$，$H=U(1)^d$ 在一般点自由作用，辛商先取 $d$ 个矩映射方程（每个降 1 个实维）、再除以 $d$ 维群（每个再降 1 个实维），共降 $2d$，故 $\dim_{\mathbb{R}}Z=2(N-d)$、$\dim_{\mathbb{C}}Z=N-d$。$F=U(1)^{N-d}$ 与 $F_{\mathbb{C}}=(\mathbb{C}^*)^{N-d}$ 即环面簇定义中的“稠密环面”。用扇（fan）的语言：$Z$ 的每个 $H$-轨道闭包对应扇中一个锥，$Z$ 的复结构由扇的粘合给出，而 $Z$ 紧致当且仅当扇是完备的（覆盖整个 $\mathbb{R}^{N-d}$）——这正是下一段紧致性判据的几何版本。]

在关于 $q_{ij}$ 的某种限制

% ==== 原文 PDF 第 57 页（印刷页 57）====
\label{op:57}
% [跨页接续: 第 56/57 页；"Under some restriction / on the qij" 跨页，按语义合译。]
下，$Z$ 是紧致的。只要对某个 $j$ 所有 $q_{ij}$ 都为正，这一点就成立。（更一般地，只要存在
正数 $m_j$ 使对所有 $i$ 都有 $\sum_j q_{ij}m_j>0$，它便成立；不过这并非实质性的推广，
因为可以通过 $U(1)^d$ 的自同构化归到前一种情形。）反之，每个紧致环面簇都可用这种方式
得到 [38,39]。
% [Tutor Note: 紧致性判据：若某个 $j$ 的所有 $q_{ij}>0$，则 $D_j=0$ 给出 $\sum_iq_{ij}|s_i|^2=r_j$（$p=0$），从而每个 $|s_i|^2\le r_j/q_{ij}$ 有界，$D$ 平坦轨迹在 $\mathbb{C}^N$ 中有界，其 $H$ 商紧致。一般条件 $\sum_j q_{ij}m_j>0$ 可通过 $U(1)^d$ 的实线性自同构（把生成元换成 $\sum_j m_jQ_j$ 之类的组合）化归为“某个 $j$ 的荷全正”的情形，故并不失去一般性。]

方程 $G=0$ 在 $Z$ 中定义超曲面 $X$。作为环面簇中的超曲面，$X$ 称为环面超曲面
（toric hypersurface）。若把 $G$ 取为横截的，意即方程
\begin{equation}
G=\frac{\partial G}{\partial S_i}=0
\tag{5.19}
\end{equation}
仅在所有 $S_i=0$ 时成立，则该超曲面在 $Z$ 的奇点之外光滑；以下只考虑这种情形。

在如上选取的 $Z$、$X$ 与适当的 $r_j$（例如恰对某个满足 $q_j>0$ 的 $j$ 取 $r_j>0$）下，
条件 $U=0$ 迫使 $s_i$ 不能全为零。于是令 $U$ 为零（并利用 $G$ 的横截性），$p$ 与 $G$
必须为零，从而在规范变换意义下经典基态的轨迹恰是超曲面 $X\subset Z$。低能理论是以该
靶空间为靶的 $\sigma$ 模型。

另一方面，通过改变 $r_j$，系统将经历许多相变。其中一些相变将涉及 $X$ 的拓扑改变，
同时保持模型（在低能下）作为 Calabi--Yau $\sigma$ 模型的解释，而改变时空的拓扑。这类
拓扑改变所涉及的局部行为将在 §5.5 中讨论。系统的其他相将是 $\sigma$ 模型与
Landau--Ginzburg 模型的混合；其中一类混合出现在本小节开头所处理的例子中。最后，若令
$r_j=-rq_j$（对某个固定的 $r$），则在规范变换意义下有唯一的经典真空；它具有 $|p|^2=r$、
$s_i=0$。此时我们得到规范 Landau--Ginzburg 模型的描述，推广了上面处理过的情形。
规范群为 $U(1)^{d-1}\times\mathbb{Z}_n$（对某个 $n$）。值得注意的是，Batyrev [40] 给出了一类% <-- 行末 % 防止跨文件拼接产生空格（s05a2 紧接续译）
% [跨页接续: 原文第 57 页/第 58 页；本句在下一页续完（"…of Landau-Ginzburg description of the
% cohomology of a general toric hypersurface"），其余部分归 s05a2-comment 分片。]
% [背景补充: 原文接着称 Batyrev 的工作给出一般环面超曲面上同调的某种 Landau--Ginzburg 描述，并推测本节所述事实为其构造提供场论背景。Batyrev 构造的组合工具是环面超曲面的牛顿多面体及其极多面体：当极多面体也是整多胞形（即两者均为反射多胞形，reflexive polytope）时，对偶的环面族给出镜像流形，其 Hodge 数可由多胞形的面数据读出。本节把“$X$ 的复结构模 = $G$ 的系数”实现为场论中的超势，把“凯勒类 = $r_j$”实现为 FI 参数，因此正好为这套组合构造提供了物理背景。]
